% ============================================================
% 3 任务与协议
% ============================================================
\section{任务与评测协议}
\label{sec:task}

\subsection{任务定义}

设真实图库 $\mathcal{G}=\{g_i\}$ 为一组真实帧，待查伪造图为 $q$。\aura{} 的目标是
在 $\mathcal{G}$ 中检索出 $q$ 的\textbf{背景来源帧} $g^\star$：
即生成 $q$ 时所沿用的那张真实图。其依据是：换脸算法保留目标视频的背景与头姿，
只替换前景人脸，故背景构成了连通 $q$ 与 $g^\star$ 的来源指纹。
训练阶段不使用任何篡改类型标签（attack-unsupervised）。

\subsection{数据与划分}

实验数据来自 FaceForensics++ C23\cite{rossler2019faceforensics} 的人脸抽帧，共 $1000$ 个真实
视频，每个真实视频对应 $5$ 种换脸算法生成的伪造视频，各抽 $8$ 帧。

伪造视频命名为 \code{X\_Y}（官方约定）：$X$ 为目标视频（提供背景与头姿，即检索答案），
$Y$ 为供体身份。于是对任意伪造帧，其背景来源真实帧为 \code{original/X} 的同序号帧，
记为 \sourceid{}。按\textbf{真实视频编号 $X$} 以 $300/100/200$ 切分为训练$/$验证$/$测试，
同一视频的帧必落在同一划分，伪造图与其源真实图也必同划分，故检索答案不会经切分泄漏。

测试协议中：\textbf{图库 $=1600$ 张测试集真实帧}，\textbf{查询 $=8000$ 张测试集伪造帧}。
随机基线的帧级 Top-1 为 $1/1600=0.0625\%$，Top-5 为 $0.3125\%$。

\subsection{检索判据}

对每条查询，按描述子距离（余弦或汉明）对全库排序，取前 $K$ 项，定义两个层级的命中：

\begin{itemize}[leftmargin=1.4em,itemsep=1pt]
\item \textbf{帧级命中}：前 $K$ 项中出现精确的 \sourceid{}（即 \code{X\_f\{帧号\}}）；
\item \textbf{视频级命中}：前 $K$ 项中出现与查询同源视频 $X$ 的真实帧。
\end{itemize}

本文统一报告帧级与视频级 Top-1$/$Top-5，并同时给出余弦与汉明两套读数。
需要说明，上述帧级判据要求帧号精确相等；由于库中每条视频只有 $8$ 帧（相隔约两秒），
它衡量的是“在相隔两秒的 $8$ 张里挑一张”。第 5.7 节另给出以\textbf{相似度区间}为命中范围的
稠密帧库协议，用于回答“能否在连续视频中定位到具体时刻”，两者的读数不可混用。

\subsection{退化鲁棒基准}
\label{sec:degrade}

真实待查的伪造图通常经历社媒转存、二次压缩与裁剪。为评测该场景，我们定义一族
\textbf{确定性}退化，\textbf{只作用于查询}，图库保持干净并只编码一次
（模拟“干净归档素材 $+$ 受损待查件”）。退化分为两族共 $12$ 档，见表~\ref{tab:levels}。

\begin{table}[t]
\centering
\caption{确定性退化协议（只作用于查询；图库保持干净）。}
\label{tab:levels}
\small
\begin{tabular}{lll}
\toprule
族 & 档位 & 说明 \\
\midrule
\multirow{6}{*}{光度}
 & \code{clean}     & 原图（对照） \\
 & \code{jpeg50}    & JPEG 质量 50 \\
 & \code{jpeg20}    & JPEG 质量 20 \\
 & \code{scale0.5}  & $0.5\times$ 下采样后插回原尺寸（丢高频） \\
 & \code{crop0.5}   & 中心裁剪 $50\%$ 后插回（丢外围背景） \\
 & \code{reencode}  & 社媒重编码链：$0.75\times \rightarrow$ JPEG70 $\rightarrow$ 插回 $\rightarrow$ JPEG60 \\
\midrule
\multirow{6}{*}{几何}
 & \code{crop0.8}   & 中心裁剪 $80\%$ 后插回 \\
 & \code{crop0.7off}& 非中心裁剪 $70\%$（固定偏移，制造整体位移） \\
 & \code{zoom1.15}  & 放大 $15\%$（中心区插回） \\
 & \code{rot5}      & 绕中心旋转 $5^\circ$ \\
 & \code{rot10}     & 绕中心旋转 $10^\circ$ \\
 & \code{shift8}    & 固定方向平移 $8$ 像素 \\
\bottomrule
\end{tabular}
\end{table}

几何族的设计动机是：光度退化只削弱高频细节，而该任务的判别信号大量存在于低频全局色调中，
故零训练的低分辨率描述子对光度档天然近乎不变；要区分不同方法，必须引入\textbf{几何}
扰动，使原始像素描述子整体错位。详细读数见第 5.4 节。
